\chapter{多元微积分}

\section{多元泰勒}

\[
f(x_0 + \Delta x, y_0 + \Delta y) 
= f(x_0, y_0) 
+ \left( \Delta x \frac{\partial}{\partial x} + \Delta y \frac{\partial}{\partial y} \right) f(x_0, y_0) 
+ \frac{1}{2!} \left( \Delta x \frac{\partial}{\partial x} + \Delta y \frac{\partial}{\partial y} \right)^2 f(x_0, y_0) 
+ \cdots 
+ \frac{1}{n!} \left( \Delta x \frac{\partial}{\partial x} + \Delta y \frac{\partial}{\partial y} \right)^n f(x_0, y_0) 
+ R_n
\]
二阶是
\[ \left( \Delta x \frac{\partial}{\partial x} + \Delta y \frac{\partial}{\partial y} \right)^2 f(x_0, y_0)
= \Delta x^2 \frac{\partial^2 f}{\partial x^2} + 2 \Delta x \Delta y \frac{\partial^2 f}{\partial x \partial y} + \Delta y^2 \frac{\partial^2 f}{\partial y^2} \]
三阶是
  \[ \left( \Delta x \frac{\partial}{\partial x} + \Delta y \frac{\partial}{\partial y} \right)^3 f(x_0, y_0)
= \Delta x^3 \frac{\partial^3 f}{\partial x^3} + 3 \Delta x^2 \Delta y \frac{\partial^3 f}{\partial x^2 \partial y} + 3 \Delta x \Delta y^2 \frac{\partial^3 f}{\partial x \partial y^2} + \Delta y^3 \frac{\partial^3 f}{\partial y^3} \]
 




\chapter{重积分}

\section{很难画出的积分区域}

\begin{example}[紫皮书6.2习题]
已知$f$在$\z[0,1\y]$上连续,证明：
\[ \int_{0}^{1}\mathrm{d} x\int_{0}^{x}\mathrm{d} y\int_{0}^{y}f(x)f(y)f(z)\mathrm{d}z = \frac{1}{3!}\z(\int_{0}^{1}f(x)\mathrm{d}x\y)^3\]
\end{example}

\begin{proof}
\begin{align*}
\text{左边}\quad
&=\int_0^1\Bigl[\int_0^x dy\int_0^y f(x)f(y)f(z)\,dz\Bigr]dx
=\int_0^1 I(x)\,dx \\
&I(x)=f(x)\int_0^x dy\int_0^y f(y)f(z)\,dz
\xlongequal{\text{交换积分次序}}
f(x)\int_0^x dz\int_z^x f(y)f(z)\,dy
\tag{*} \\
&\xlongequal{\text{交换 }y,z\text{ 次序}}
f(x)\int_0^x dy\int_y^x f(z)f(y)\,dz \\
&\xlongequal{\text{与原式 }(*)\text{ 相加}}
\frac12 f(x)\int_0^x dy\int_0^x f(y)f(z)\,dz \\
&=\frac12 f(x)\int_0^x f(y)\,dy\int_0^x f(z)\,dz \\
\text{所以}\quad
&\int_0^1 I(x)\,dx
=\int_0^1 \frac12 f(x)\int_0^x f(y)\,dy\int_0^x f(z)\,dz\,dx \\
&=\frac12\int_0^1
\underbrace{f(x)\,dx}_{\textcolor{red}{\text{凑}d\bigl[\int_0^x f(t)\,dt\bigr]}}
\Bigl[\int_0^x f(t)\,dt\Bigr]^2 \\
&=\frac12\int_0^1 \Bigl[\int_0^x f(t)\,dt\Bigr]^2
\,d\Bigl[\int_0^x f(t)\,dt\Bigr] \\
&=\frac12\cdot\frac13
\Bigl[\int_0^x f(t)\,dt\Bigr]^3\Big|_0^1 \\
&=\frac16\Bigl(\int_0^1 f(t)\,dt\Bigr)^3.
\end{align*}
\end{proof}

\begin{remark}
这是去年我做时，没做出，以为是什么几何对称性，
（\(V=\)棱锥\(=\dfrac16V_{\text{正方体}}\)）

但今天我用纯积分方法做，感觉其实更好

这种方法对 \(3\) 次以上也可以做
\end{remark}

\begin{example}[\;17届CMC]
设 \(\Omega\) 是由平面
$\Pi_1 : x = 0, \Pi_2 : x = y, \Pi_3 : y = z, \Pi_4 : z = 1, \Pi_5 : z = 2$

围成的空间区域，求三重积分$\displaystyle \iiint_{\Omega} \frac{\mathrm{d} x \, \mathrm{d} y\, \mathrm{d} z}{\sqrt{y^2 + z^2}}.$

\end{example}

\begin{note}
这一题可以使用画正方体的方法，来作出积分区域。
\end{note}



\chapter{曲面积分}

\section{算子法}

\section{凑微分}

\section{利用物理结论}


\chapter{级数}

\begin{example}
证明$\displaystyle \int_{0}^{+\infty} \frac{1}{1+x^4\sin^2 x} \,\mathrm{d}x$收敛，$\displaystyle \int_{0}^{+\infty} \frac{1}{1+x^2\sin^2 x}\,\mathrm{d}x $发散
\end{example}

\begin{proof}
\begin{align*}
\text{原级数}
&=\sum\limits_{n=0}^{\infty}\int_{n\pi}^{(n+1)\pi}\frac{\mathrm{d}x}{1+(x\pi)^4\sin^2 x}\\
& \le \int_{0}^{\pi} \frac{\mathrm{d}x}{1+x^4\sin^2 x}+\sum\limits_{n=1}^{\infty}\int_{0}^{\pi}\frac{\mathrm{d}x}{1+(n\pi)^4\sin^2 x} \\
&=I +  2\sum\limits_{n=1}^{\infty}\int_{0}^{\frac{\pi}{2}}\frac{\mathrm{d}x}{1+(n\pi)^4\sin^2 x}\\
&=I + 2\sum\limits_{n=1}^{\infty}\int_{0}^{\frac{\pi}{2}}\frac{\mathrm{d}\tan x}{1+[(n\pi)^4+1]\tan^2 x}\\
&=I + 2 \sum\limits_{n=1}^{\infty} \frac{1}{\sqrt{(n\pi)^4 +1}} \arctan (\sqrt{(n\pi)^4+1} \tan x)\\
&\le I+2\sum\limits_{n=1}^{\infty} \frac{1}{\sqrt{(n\pi)^4 +1}} \frac{\pi}{2}
\end{align*}

而 $\displaystyle \sum\limits_{n=1}^{\infty} \frac{1}{\sqrt{(n\pi)^4 +1}}$显然收敛（用比较判别法极限形式）
\end{proof}